% =============================================================================
% book.tex — Pandoc template for "Agéntico por Diseño, Tomo I"
% =============================================================================
\documentclass[$if(classoption)$$for(classoption)$$classoption$$sep$,$endfor$$endif$]{paradigma-agentico}

% Load additional style packages (found via TEXINPUTS)
\input{pa-colors.sty}
\input{pa-callouts.sty}
\input{pa-tables.sty}
\input{pa-codeblocks.sty}
\input{pa-diagrams.sty}
\input{pa-typography.sty}

% Metadata
\title{Agéntico por Diseño, Tomo I: Tecnologías de la Información}
\author{Karim Touma}
\date{$book.date$}

% Language is set via babel in the .cls file

\begin{document}

% =============================================================================
% FRONTMATTER
% =============================================================================
\frontmatter
\pagestyle{empty}

% --- Half title page ---
\vspace*{\fill}
\begin{center}
  {\fontsize{28}{34}\selectfont\sffamily\bfseries\color{pa-primary}Agéntico por Diseño}\\[12pt]
  {\fontsize{16}{20}\selectfont\sffamily\color{pa-midgray}Tomo I}
\end{center}
\vspace*{\fill}
\clearpage

% --- Full title page ---
\vspace*{0.18\textheight}
\begin{center}
  {\fontsize{13}{16}\selectfont\sffamily\color{pa-midgray}\MakeUppercase{Agéntico}}\\[0.8em]
  {\fontsize{34}{40}\selectfont\sffamily\bfseries\color{pa-primary}por Diseño}\\[1.2em]
  {\color{pa-primary}\rule{0.45\textwidth}{2pt}}\\[1.2em]
  {\fontsize{14}{18}\selectfont\sffamily\bfseries\color{pa-secondary}TOMO I}\\[0.8em]
  {\fontsize{18}{22}\selectfont\sffamily\color{pa-darktext}Tecnologías de la Información}\\[2.5em]
  {\fontsize{13}{16}\selectfont\sffamily\color{pa-midgray}Karim Touma}\\[0.6em]
  {\fontsize{10}{13}\selectfont\sffamily\color{pa-midgray}\href{mailto:karim@touma.io}{karim@touma.io}\enspace\textbar\enspace\href{https://karim.touma.io}{karim.touma.io}}\\[0.8em]
  {\fontsize{11}{14}\selectfont\sffamily\color{pa-midgray}$book.date$}
\end{center}
\vspace*{\fill}
\clearpage

% --- Copyright page ---
\vspace*{\fill}
\begin{flushleft}
  \small\sffamily\color{pa-darktext}
  \textbf{Agéntico por Diseño, Tomo I: Tecnologías de la Información}\\[12pt]
  Karim Touma, $book.date$\\[6pt]
  \href{mailto:karim@touma.io}{karim@touma.io}\enspace\textbar\enspace\href{https://karim.touma.io}{karim.touma.io}\\[18pt]
  Esta obra se libera al dominio p\'{u}blico bajo la licencia \textbf{Unlicense}.\\[6pt]
  Eres libre de copiarla, modificarla, publicarla y distribuirla, con o sin
  fines comerciales, sin pedir permiso ni dar atribuci\'{o}n.\\[6pt]
  Texto completo de la licencia y c\'{o}digo fuente del libro:\\[2pt]
  \href{https://github.com/karimtouma/agentico}{github.com/karimtouma/agentico}\\[18pt]
  \textbf{Primera edición}, $book.date$\\[6pt]
  Diseñado y compuesto con \LaTeX{}
\end{flushleft}
\clearpage

% --- Cómo Usar Este Libro ---
\pagestyle{paradigma-plain}
\vspace*{0.5cm}
{\centering\fontsize{20}{24}\selectfont\sffamily\bfseries\color{pa-primary}Cómo Usar Este Libro\par}
\vspace{1cm}

\noindent\color{pa-darktext}
Este no es un libro para leer de principio a fin y guardar en el estante.
Es una \textbf{herramienta de consulta} diseñada para que la uses en tu trabajo diario:
anota en los márgenes, subraya, marca páginas, y vuelve a las secciones que necesites.

\vspace{1em}

\noindent{\sffamily\bfseries\color{pa-primary}Estructura del libro}

\vspace{0.5em}
\noindent El contenido se organiza en seis partes que puedes abordar en cualquier orden:

\vspace{0.5em}
\begin{description}[leftmargin=0.5em,itemsep=4pt,font=\sffamily\bfseries\color{pa-secondary}]
  \item[Parte I - Contexto Estratégico] Por qué la IA agéntica importa \emph{ahora} (Caps.\ 0--4)
  \item[Parte II - Sesgos y Evidencia] Los errores cognitivos que debes evitar y qué dice la evidencia por industria (Caps.\ 5--6)
  \item[Parte III - La Tecnología] Evolución técnica y ecosistema de herramientas (Caps.\ 7--8)
  \item[Parte IV - Impacto en el Negocio] ROI, TCO, y qué pasa cuando falla (Caps.\ 9--10)
  \item[Parte V - Liderazgo y Estrategia] Cómo liderar equipos y diseñar la adopción (Caps.\ 11--12)
  \item[Parte VI - Gobernanza y Futuro] Riesgos, regulación, y visión 2030 (Caps.\ 13--14)
\end{description}

\vspace{1em}

\noindent{\sffamily\bfseries\color{pa-primary}Si necesitas\ldots}

\vspace{0.5em}
\begin{tblr}{
  colspec = {X[1.6,l] X[1,l]},
  row{1} = {font=\sffamily\bfseries\color{white}, bg=pa-primary},
  row{even} = {bg=pa-tablealt},
  hlines = {pa-tablehead},
  rowsep = 4pt,
  colsep = 8pt,
}
  Tu situación & Ve a\ldots \\
  Convencer al board de invertir en IA & Caps.\ 1, 9, App.\ B \\
  Evaluar herramientas agénticas & Caps.\ 8, 11, App.\ B \\
  Armar el business case / ROI & Caps.\ 9, 12, App.\ C \\
  Liderar el cambio cultural & Caps.\ 5, 11, 13 \\
  Diseñar gobernanza y compliance & Cap.\ 13 \\
  Entender qué es IA agéntica & Caps.\ 3, 4, 7 \\
  Identificar quick wins por industria & Cap.\ 6 \\
  Ver qué puede salir mal & Cap.\ 10 \\
\end{tblr}

\clearpage

\noindent{\sffamily\bfseries\color{pa-primary}Lectura recomendada por rol}

\vspace{0.5em}
\begin{tblr}{
  colspec = {X[0.8,l] X[1.4,l] X[0.5,c]},
  row{1} = {font=\sffamily\bfseries\color{white}, bg=pa-primary},
  row{even} = {bg=pa-tablealt},
  hlines = {pa-tablehead},
  rowsep = 4pt,
  colsep = 8pt,
}
  Rol & Capítulos clave & Tiempo \\
  CTO / VP Engineering & 1, 3, 7, 8, 11, 13 & 3--4 hrs \\
  CFO / Finance & 1, 9, 10, 12, App.\ B & 2--3 hrs \\
  Tech Lead / Architect & 3, 4, 7, 8, 11, App.\ E & 3--4 hrs \\
  CHRO / People & 5, 11, 13 & 2--3 hrs \\
  CEO / Board Member & Brief, 1, 9, 14 & 1--2 hrs \\
\end{tblr}

\vspace{1.5em}

\noindent{\sffamily\bfseries\color{pa-primary}Leyenda de elementos}

\vspace{0.5em}
\noindent A lo largo del libro encontrarás recuadros con información especial:

\vspace{0.5em}
\begin{description}[leftmargin=0.5em,itemsep=3pt,font=\sffamily\small\color{pa-secondary}]
  \item[\faClipboardList\enspace Resumen Ejecutivo] Los puntos clave del capítulo en 60 segundos
  \item[\faUsers\enspace Para Tu Reunión de Liderazgo] Ejercicios para aplicar en tu próxima reunión
  \item[\faChartBar\enspace Dato Verificado] Estadísticas con fuente, metodología y limitaciones
  \item[\faExclamationTriangle\enspace Errores Comunes] Trampas frecuentes que líderes deben evitar
  \item[\faCheckSquare\enspace Checklist] Listas de verificación para implementación
  \item[\faLightbulb\enspace Ejemplo Práctico] Aplicación concreta de un concepto
\end{description}

\vspace{1em}

\noindent{\sffamily\bfseries\color{pa-primary}El margen es tuyo}

\vspace{0.5em}
\noindent Los márgenes amplios están diseñados para que escribas. Anota ideas, marca lo que
aplica a tu contexto, conecta secciones entre sí. Un libro de consulta se vuelve más valioso
cada vez que lo usas.

\clearpage

% --- Table of Contents ---
{
  \hypersetup{linkcolor=pa-darktext}
  \setcounter{tocdepth}{1}
  \tableofcontents*
}
\clearpage

% =============================================================================
% MAINMATTER
% =============================================================================
\mainmatter
\pagestyle{paradigma-page}
\markboth{}{}

% Part dividers are inserted automatically by part-dividers.lua filter
% based on chapter numbering (Parts I-VI)

$body$

% =============================================================================
% BACKMATTER
% =============================================================================
\backmatter

% --- Índice Analítico ---
\clearpage
\pagestyle{paradigma-plain}
\printindex

% --- Sobre el Autor ---
\clearpage
\pagestyle{paradigma-plain}
\vspace*{0.3cm}
{\centering\fontsize{22}{26}\selectfont\sffamily\bfseries\color{pa-primary}Sobre el Autor\par}
\vspace{0.8cm}

\noindent\color{pa-darktext}
\textbf{Karim Touma} ha construido su carrera en la intersecci\'{o}n entre
estrategia ejecutiva y ejecuci\'{o}n t\'{e}cnica: el l\'{i}der que toma decisiones
en la mesa directiva y tambi\'{e}n entiende qu\'{e} sucede dentro del modelo.

\vspace{0.6em}
\noindent Como investigador en la Universidad de Santiago de Chile, dise\~{n}\'{o}
modelos de data mining cuando el machine learning a\'{u}n era territorio acad\'{e}mico.
En Equifax hizo I+D a escala continental. En Telef\'{o}nica Chile construy\'{o} la
primera capacidad de Big Data Analytics de una telco en Latinoam\'{e}rica. En
Falabella posicion\'{o} la anal\'{i}tica como capacidad estrat\'{e}gica. En LATAM
Airlines despleg\'{o} personalizaci\'{o}n en tiempo real para millones de pasajeros.
Fund\'{o} una startup de pricing basado en datos. Lider\'{o} equipos de m\'{a}s de
cien ingenieros en seis pa\'{i}ses.

\vspace{0.6em}
\noindent A lo largo del camino, ha analizado din\'{a}micas de comportamiento de
m\'{a}s de 400 millones de clientes - telecomunicaciones, retail, aviaci\'{o}n,
cr\'{e}dito - y en cada industria observ\'{o} el mismo patr\'{o}n: la tecnolog\'{i}a
sin redise\~{n}o produce mejoras incrementales; la tecnolog\'{i}a con redise\~{n}o
produce transformaciones.

\vspace{0.6em}
\noindent Hoy, como Group Chief Digital \& Data Officer de Grupo DEACERO, lidera
la transformaci\'{o}n digital de una de las mayores manufactureras de acero de
M\'{e}xico. Arquitecta personalmente la plataforma de GenAI del grupo, construye
sistemas RAG, hace fine-tuning de LLMs y los despliega a escala en
producci\'{o}n. Es la quinta industria radicalmente distinta en la que aplica la
misma convicci\'{o}n: que la IA solo genera valor cuando se dise\~{n}a alrededor
de ella, no cuando se agrega encima de lo que ya existe.

\vspace{0.6em}
\noindent Es alumni del Stanford Graduate School of Business y graduado en Ciencias
de la Computaci\'{o}n de la Universidad de Santiago de Chile.

\vspace{0.6em}
\noindent Este libro destila dos d\'{e}cadas navegando entre la mesa directiva y el
terminal. No es la perspectiva del acad\'{e}mico ni la del consultor externo. Es
la del operador que ha apostado su carrera a una idea simple: que la forma de
capturar valor de cada revoluci\'{o}n tecnol\'{o}gica no es adoptar herramientas,
sino \textbf{redise\~{n}ar organizaciones}.

\vspace{1.2em}
\begin{center}
\sffamily\small
\faEnvelope\enspace\href{mailto:karim@touma.io}{karim@touma.io}\qquad
\faGlobe\enspace\href{https://karim.touma.io}{karim.touma.io}\qquad
\faLinkedin\enspace\href{https://www.linkedin.com/in/katouma}{linkedin.com/in/katouma}
\end{center}

\end{document}
